The \wikilink{s28--b005}{prime-ideal correspondence} is related to the
following calculation for arbitrary ideals. Let $A$ be commutative and
unital, and let $S$ contain $1$ and be multiplicatively closed.
Write $j:A\to S^{-1}A$.

\begin{definition}\label{ideal-extension-contraction}
The \emph{extension} of an ideal $I\subseteq A$ is $S^{-1}I$, the ideal
generated by $j(I)$. The \emph{contraction} of an ideal $J\subseteq S^{-1}A$
is $j^{-1}J$. The \emph{$S$-saturation} of $I$ is
$I^{\mathrm{sat}}=\{a:sa\in I\text{ for some }s\in S\}$.
An ideal is \emph{$S$-saturated} if $I=I^{\mathrm{sat}}$.
\end{definition}
\begin{proposition}\label{localized-ideals}
Every ideal $J$ of $S^{-1}A$ equals $S^{-1}(j^{-1}J)$, and
$j^{-1}(S^{-1}I)=I^{\mathrm{sat}}$ for every ideal $I$ of $A$.
Thus localization gives a bijection with the $S$-saturated ideals of $A$.
\end{proposition}
\begin{proof}\label{localized-ideals-proof}
If $a/s\in J$, multiplying by $s/1$ gives $a/1\in J$.
Conversely, $a\in j^{-1}J$ implies $a/s\in J$ because $s/1$ is a unit.
This proves the first equality. Elements of $S^{-1}I$ have the form
$b/s$ with $b\in I$. If $a/1=b/s$, there is $t\in S$ with
$t(sa-b)=0$, so $(ts)a\in I$. Conversely, $sa\in I$ implies
$a/1=(sa)/s\in S^{-1}I$. Contractions are saturated since $j(s)$
is a unit; the equalities now establish the bijection.
\end{proof}
\begin{proposition}\label{adjoin-inverse}
For $f\in A$, there is a natural isomorphism
$A[y]/(fy-1)\cong A[f^{-1}]$, the localization at $\{1,f,f^2,\ldots\}$.
\end{proposition}
\begin{proof}\label{adjoin-inverse-proof}
An $A$-algebra map from the quotient to a commutative $A$-algebra $B$
is a choice of $b$ with $fb=1$. It exists exactly when $f$ becomes a
unit, and is then unique. This is the
\wikilink{s26--b006}{localization universal property}. The maps sending
$y$ to $f^{-1}$ and $f^{-1}$ to the class of $y$ are inverse.
\end{proof}
\begin{remark}\label{localized-ideals-scope}
There is no bijection with all ideals of $A$. Inverting $2$ in
$\mathbb Z$ makes $(2)$ extend to the whole ring, whose contraction
is $\mathbb Z$, the saturation of $(2)$. The propositions allow zero
divisors. If $0\in S$, the localized ring is zero and the equalities
remain valid.
\end{remark}
